% ── SLIDE 4: Bottle Neck ──────────────────────────────────────────────
\begin{frame}{YOLOv11: Bottle Neck Block}
  \begin{columns}[c]
    \begin{column}{0.40\textwidth}
      \centering
      \begin{tikzpicture}[
          outer/.style={
              draw=utecblue!85,
              fill=utecblue!5,
              rounded corners=15pt,
              line width=1.2pt,
              minimum width=3.35cm,
              minimum height=5.25cm
            },
          layer/.style={
              draw=utecblue!85,
              fill=white,
              rounded corners=7pt,
              line width=1.0pt,
              minimum width=2.35cm,
              minimum height=0.62cm,
              font=\small\bfseries,
              align=center
            },
          sum/.style={
              circle,
              draw=utecblue!85,
              fill=white,
              line width=1.0pt,
              minimum size=0.62cm,
              font=\large\bfseries
            },
          arr/.style={-{Stealth[length=5pt]}, utecblue!80, line width=1.0pt},
          shortcut/.style={-{Stealth[length=5pt]}, utecblue!80, line width=1.1pt}
        ]
        \node[outer] at (0,0) {};
        \node[font=\Large\bfseries, text=utecblue!90] at (0,2.13) {Bottle Neck};

        \node[layer] (cbs1) at (0,1.08) {CBS};
        \node[layer] (cbs2) at (0,-0.18) {CBS};
        \node[sum] (add) at (0,-1.45) {$+$};

        \coordinate (input) at (0,1.75);
        \draw[arr] (input) -- (cbs1.north);
        \draw[arr] (cbs1) -- (cbs2);
        \draw[arr] (cbs2) -- (add);

        \draw[shortcut] (input) -- ++(-1.45,0) |- (add.west);
        \node[font=\scriptsize, text=utecblue!80] at (-1.06,1.62) {shortcut};
      \end{tikzpicture}
    \end{column}
    \begin{column}{0.56\textwidth}
      \textbf{Bottleneck idea}
      \begin{itemize}
        \item Two CBS layers transform the feature map with low extra cost.
        \item The shortcut adds the input back to the output:
      \end{itemize}
      \[
        \mathbf{y} = F(\mathbf{x}) + \mathbf{x}
      \]
      \begin{itemize}
        \item Same residual principle as a ResNet block: preserve useful input information and improve gradient flow.
        \item CBS means \textbf{Conv + BatchNorm + SiLU}; convolution extracts features, BatchNorm stabilizes training, and SiLU adds nonlinearity.
      \end{itemize}
    \end{column}
  \end{columns}
\end{frame}

